\documentclass[10pt,a4paper]{amsart}
\usepackage[margin=19mm]{geometry}
\usepackage{amsmath,amssymb,mathtools}
\usepackage[scr=rsfs,cal=euler]{mathalfa}
\usepackage{xfrac}
\usepackage[unicode,hidelinks,bookmarks=false]{hyperref}
\newcommand{\ie}{i.e.\ }
\newcommand{\cf}{cf.\ }
\newcommand{\resp}{respectively\ }
\newcommand{\Subsection}{\subsection}
\providecommand{\nobreakdash}{\nobreak\hbox{-}}
\setlength{\emergencystretch}{2em}
\multlinegap0pt
\usepackage{amssymb}
\usepackage{bm}
\usepackage{enumerate}
\usepackage{xcolor}
\newtheorem{proposition}{Proposition}
\title{Understanding and Resolving Singularities in 3D Dirichlet Boundary Problems}
\author{David Levin}
\date{Source: arXiv:2603.09926v1 (CC BY 4.0)}
\begin{document}
\maketitle

\noindent\emph{Source and licence.} Understanding and Resolving Singularities in 3D Dirichlet Boundary Problems. Version arXiv:2603.09926v1 (CC BY 4.0), distributed by arXiv under the Creative Commons Attribution 4.0 International licence, http://creativecommons.org/licenses/by/4.0/, as recorded in the arXiv licence metadata for this exact version and shown on the abstract page https://arxiv.org/abs/2603.09926v1. Full licence notice: \texttt{licenses/arxiv-2603.09926-CC-BY-4.0.txt}.\par\smallskip

\noindent\textit{Editorial scope.} Counted unit: the article's proposition proposition (source0.tex lines 683--697). The article's complete original proof of it is delivered with it (source0.tex lines 699--726, one \texttt{proof} environment of 28 lines). No mathematics is written, repaired or re-derived: the statement and the proof are the article's own lines, in the article's own notation, and no retained line is substituted or re-flowed.  Retained uncounted as the source-local context the proof needs: lines 258--264; lines 289--297; lines 554--557; lines 589--592.  Nothing was omitted from any retained span, so the package carries no omission record.  No retained line contains a citation, so the wrapper carries no bibliography and no bibliography entry.  No retained line contains a figure environment, an image-inclusion command or drawing code, so no image is delivered and the package declares no asset dependency.  Editorial formatting only, in addition: an AMS \texttt{amsart} preamble that restates the article's own declarations and macros used by the retained lines, namely the environments \texttt{proposition} on separate counters, together with the standard packages those lines need.  The retained environments print in this document as Proposition 1 in order of appearance, whereas the article prints the same items with the numbering of its own full document.  The complete native source of the article as distributed in the arXiv e-print for this exact version is bundled unchanged as \texttt{sources/196101/source0.tex}.
\begin{equation}\label{Gcube}
G^\square(x,y)
= \sum_{n\in\mathbb{Z}^3}
\sum_{m\in\{0,1\}^3}
(-1)^{m_1+m_2+m_3}\;
\Phi\!\left(x - y_{n,m}\right),
\end{equation}

\begin{equation}\label{G27}
G^\square_{\text{27}}(x,y)
= \sum_{i,j,k=-1}^{1} 
  \sum_{m_1,m_2,m_3=0}^{1} 
  (-1)^{m_1+m_2+m_3}\;
  \Phi\Bigl(
    x - y^{(i,j,k)}_{m_1,m_2,m_3}
  \Bigr),
\end{equation}

\begin{equation}
G^\square(x,y) = G^\square_{27}(x,y) + R^\square_{27}(x,y),
\label{SRcube}
\end{equation}

\begin{equation}\label{H_R}
H_R(x)\equiv \int_{\partial\Omega}
    \frac{\partial R}{\partial n_y}(x,y)f(y)\, ds_y .
\end{equation}

\begin{proposition}
Let $\Omega=[0,1]^3$ and consider the decomposition
\eqref{SRcube}, where $G_{\square}$ is given by the lattice
representation \eqref{Gcube} and $G_{\square}^{27}$ by the
truncated sum \eqref{G27}.  
Then, for every $y\in\Omega$, the remainder
$R_{\square}^{27}(\cdot,y)$ is harmonic in the enlarged domain
$(-1,2)^3$.  Consequently, the function
\[
H_R(x)=\int_{\partial\Omega}
\frac{\partial R_{\square}^{27}}{\partial n_y}(x,y)\,
f(y)\,dS_y,
\]
defined in \eqref{H_R}, is harmonic in $(-1,2)^3$.
\end{proposition}

\begin{proof}
By \eqref{Gcube}, the Green’s function is represented as a
lattice sum of free-space source functions
$\Phi(x-y_{n,m})$, while \eqref{G27} retains only those
terms with $n\in\{-1,0,1\}^3$.
Hence, in the decomposition \eqref{SRcube}, the remainder
$R_{\square}^{27}$ consists precisely of those image terms
with $n\notin\{-1,0,1\}^3$.
For any $y\in[0,1]^3$, such image points $y_{n,m}$ lie
outside the box $(-1,2)^3$.
Therefore, each term
$x\mapsto\Phi(x-y_{n,m})$ is harmonic throughout
$(-1,2)^3$.
On compact subsets of $(-1,2)^3$, the series defining
$R_{\square}^{27}(\cdot,y)$ admits termwise differentiation
with respect to $x$ (for instance, via uniform convergence
of the second-derivative series),
and thus $\Delta_x R_{\square}^{27}(x,y)=0$
in $(-1,2)^3$.
Differentiation with respect to $y$ does not affect
harmonicity in $x$, and hence
$\partial R_{\square}^{27}/\partial n_y(\cdot,y)$
is harmonic in $x$ on $(-1,2)^3$.
Finally, the representation \eqref{H_R}
shows that $H_R$ is an integral of harmonic functions
with respect to $y$, and is therefore harmonic
in $(-1,2)^3$.
\end{proof}

\end{document}
